% ============================================================
% Controlled modality analysis — FINAL LaTeX table (v2)
% Paste into Section 4.12 of DS_RangeNet_MDPI.tex
% Location: after the existing ablation discussion
% ============================================================

\begin{table}[H]
\centering
\caption{Controlled modality analysis on the UBPC-9 validation set.
Inference-time perturbation experiments use the same DS-RangeNet
checkpoint; the corresponding input channels are zero-masked or
perturbed while all network weights remain unchanged. Separately
trained modality-ablation models are identified explicitly. Each
perturbation is repeated three times with different random seeds;
we report mean $\pm$ sample standard deviation over the three trials.}
\label{tab:modality_controls}
\begin{tabular}{lcc}
\toprule
\textbf{Control} & \textbf{Input condition} & \textbf{mIoU (\%)} \\
\midrule
\multicolumn{3}{l}{\textit{Separately trained single-/early-fusion models}} \\
\quad Geometry-only model   & Geometry stream only (reflectance removed) & $58.3\pm0.35$ \\
\quad Reflectance-only model & Reflectance stream only (geometry removed) & $56.8\pm0.26$ \\
\quad Early fusion model    & Aligned 16-channel concatenation (separately trained) & $67.9\pm0.35$ \\
\midrule
\multicolumn{3}{l}{\textit{Inference-time perturbations (same DS-RangeNet checkpoint)}} \\
\quad Intensity missing     & Reflectance channels zero-masked in the full model & $65.4\pm0.26$ \\
\quad Intensity corrupted   & $I'=\operatorname{clip}(1.2I+\epsilon,\,0,\,1)$, $\epsilon\sim\mathcal{N}(0,0.05^2)$ & $66.1\pm0.26$ \\
\quad Geometry sparse       & Random point removal (30\%); voxel-PCA and curvature \emph{recomputed} & $67.3\pm0.26$ \\
\quad Cross-frame mismatch  & Geometry from $t$, reflectance from $t{+}\delta$ within the \emph{same} sequence; $\delta\in\{1,5,10,20\}$ chosen randomly & $64.7\pm0.32$ \\
\midrule
\textbf{DS-RangeNet}       & Correctly aligned modalities (baseline) & $\mathbf{73.2}$ \\
\bottomrule
\end{tabular}
\end{table}

% ============================================================
% Implementation details (paste as a paragraph after the table)
% ============================================================
\textbf{Implementation details.}
The intensity values are normalized to $[0,1]$; therefore
$\sigma=0.05$ denotes Gaussian noise with a standard deviation equal to
5\% of the normalized intensity range. The intensity shift is applied as
$I'=\operatorname{clip}(1.2I+\epsilon,\,0,\,1)$ with clipping to
$[0,1]$. For cross-frame mismatch, the geometry scan comes from frame
$t$ and the reflectance scan is sampled from frame $t+\delta$ within the
\emph{same} sequence (offset $\delta$ randomly chosen from
$\{1,5,10,20\}$), ensuring that the experiment isolates spatial
misalignment rather than scene change. For geometry sparse, voxel-PCA
descriptors and intensity curvature are \emph{recomputed} on the
remaining 70\% of points; the descriptors are not frozen from the
original point cloud. Each perturbation is repeated three times with
different random seeds, and we report mean $\pm$ sample standard
deviation. Under the present training and masking protocol, the
geometry-only condition performs slightly better than the
reflectance-only condition, suggesting that geometry provides a
stronger standalone cue in this dataset. Among the two-modality
perturbation controls, cross-frame mismatch causes the largest
degradation ($64.7\%\pm0.25$), confirming the importance of
spatially aligned coupling between geometry and reflectance.
